A carta de intenções é uma forma textual proposta para se pensarem práticas científicas politeístas orientadas pela e para a comunicação. Seu texto apresenta proposições científicas a serem elaboradas, articuladas ou comentadas em seu meta-texto como notas de rodapé iteráveis. Apresenta-se, assim, como um convite para que as notas sejam refeitas enquanto re-localizações do texto, com avaliações, novas articulações, explicações, comentários, expansões, etc; e também para que novas notas sejam inseridas em trechos até então não comentados.